% Part: turing-machines
% Chapter: undecidability
% Section: representing-tms

\documentclass[../../../include/open-logic-section]{subfiles}

\begin{document}

\olfileid{tur}{und}{rep}
\olsection{ట్యూరింగ్ యంత్రాలకు ప్రాతినిధ్యం}

\begin{explain}
ట్యూరింగ్ యంత్రాలకూ వాటి ప్రవర్తనకూ మొదటిస్థాయి తర్కంలోని \tetoken{ఒక వాక్యంతో}{!!a{sentence}}
ప్రాతినిధ్యం వహించాలంటే, తగిన భాషను నిర్వచించాలి. ఆ భాషలో రెండు భాగాలు
ఉంటాయి: యంత్రపు స్థితివిన్యాసాలను వర్ణించే \tetoken{విధేయ సంకేతాలు}{!!{predicate}s}, అమలు అంచెలు
(అంటే ``క్షణాలు''), టేపుపై స్థానాలకు సంఖ్యలు ఇచ్చే వ్యక్తీకరణలు.

రెండు రకాల \tetoken{విధేయ సంకేతాలను}{!!{predicate}s} ప్రవేశపెడతాం; రెండూ ద్విస్థానికాలు: ప్రతి
స్థితి~$q$కు \tetoken{ఒక విధేయ సంకేతం}{!!a{predicate}}~$\Obj Q_q$, ప్రతి టేపు సంకేతం~$\sigma$కు
\tetoken{ఒక విధేయ సంకేతం}{!!a{predicate}}~$\Obj S_\sigma$. మొదటి రకం $M$ స్థితినీ దాని టేపు శీర్షం
స్థానాన్నీ వర్ణించడానికి, రెండవ రకం టేపులోని విషయాన్ని వర్ణించడానికి
ఉపయోగపడతాయి.

టేపు శీర్షం స్థానాలనూ అమలైన అంచెల సంఖ్యనూ వ్యక్తం చేయడానికి, సంఖ్యలను
వ్యక్తం చేసే పద్ధతి కావాలి. దీనికి \tetoken{ఒక స్థిరసంకేతం}{!!a{constant}}~$\Obj 0$, అలాగే
ఉత్తరవర్తి ప్రమేయమైన $1$-స్థాన ప్రమేయ సంకేతం~$\prime$ను ఉపయోగిస్తాం.
ఆచారం ప్రకారం దాన్ని దాని ఆర్గ్యుమెంట్ \emph{తరువాత} వ్రాస్తాం (కుండలీకరణాలను
వదిలేస్తాం).

ప్రతి సంఖ్య $n$కు $\Lang L_M$లో దానికి ప్రాతినిధ్యం వహించే నియత పదం
$\num{n}$ ఉంటుంది; ఇది $n$కు \emph{సంఖ్యాపదం}. $\num{0}$ అంటే $\Obj 0$,
$\num{1}$ అంటే $\Obj 0'$, $\num{2}$ అంటే $\Obj 0''$, ఇలాగే కొనసాగుతుంది.
మరింత ఆచారబద్ధంగా:
\begin{align*}
\num{0} & = \Obj 0 \\
\num{n+1} &= \num{n}'
\end{align*}
$\num{0}$ అనే పదం, అంటే $\Obj 0$, టేపుపై అత్యంత ఎడమ స్థానాన్నీ, మొదటి
అమలు అంచెకు ముందున్న సమయాన్నీ (ప్రారంభ స్థితివిన్యాసాన్ని) సూచిస్తుంది.
$\num{1}$ అనే పదం, అంటే $\Obj 0'$, అత్యంత ఎడమ గడికి కుడివైపున్న గడినీ,
మొదటి అమలు అంచె తరువాతి సమయాన్నీ సూచిస్తుంది; ఇలాగే కొనసాగుతుంది.

టేపు స్థానాల క్రమాన్ని (అప్పుడు దాని అర్థం ``ఎడమవైపున'') మరియు అమలు అంచెల
క్రమాన్ని (అప్పుడు దాని అర్థం ``ముందు'') వ్యక్తం చేయడానికి \tetoken{ఒక విధేయ సంకేతం}{!!a{predicate}}~$<$ను
కూడా ప్రవేశపెడతాం.

భాష సిద్ధమైన తరువాత $!T(M,w)$ యొక్క ``స్వీకృతాలను,'' అంటే $M$ను ఇన్‌పుట్~$w$పై
నడిపినప్పుడు దాని ప్రవర్తనను కలిపి వర్ణించే \tetoken{వాక్యాలను}{!!{sentence}s}, జాబితా చేస్తాం.
$\Obj 0$, $\prime$, $<$లపై షరతులు విధించే \tetoken{వాక్యాలు}{!!{sentence}s}, ఇన్‌పుట్
స్థితివిన్యాసాన్ని వర్ణించే \tetoken{వాక్యాలు}{!!{sentence}s}, ఒక నిర్దిష్ట నిర్దేశాన్ని $M$
అమలు చేసిన తరువాత దాని స్థితివిన్యాసాన్ని వర్ణించే \tetoken{వాక్యాలు}{!!{sentence}s} ఉంటాయి.
\end{explain}

\begin{defn}
  \ollabel{defn:tm-descr}
ట్యూరింగ్ యంత్రం $M = \tuple{Q, \Sigma, q_0, \delta}$ ఇచ్చినప్పుడు,
భాష~$\Lang L_M$లో ఇవి ఉంటాయి:
\begin{enumerate}
\item ప్రతి స్థితి~$q \in Q$కు ఒక ద్విస్థాన \tetoken{విధేయ సంకేతం}{!!{predicate}}
  $\Obj Q_q(x, y)$. అంతర్బోధాత్మకంగా, $\Obj Q_q(\num{m}, \num{n})$ అంటే
  ``$n$ అంచెల తరువాత, $M$ స్థితి~$q$లో ఉండి $m$వ గడిని చదువుతోంది.''
\item ప్రతి సంకేతం~$\sigma\in \Sigma$కు ఒక ద్విస్థాన \tetoken{విధేయ సంకేతం}{!!{predicate}}
  $\Obj S_\sigma(x, y)$. అంతర్బోధాత్మకంగా,
  $\Obj S_\sigma(\num{m}, \num{n})$ అంటే ``$n$ అంచెల తరువాత, $m$వ గడిలో
  సంకేతం~$\sigma$ ఉంది.''
\item ఒక \tetoken{స్థిరసంకేతం}{!!{constant}} $\Obj 0$
\item ఒక ఏకస్థాన \tetoken{ప్రమేయ సంకేతం}{!!{function}} $\prime$
\item ఒక ద్విస్థాన \tetoken{విధేయ సంకేతం}{!!{predicate}} $<$
\end{enumerate}
\end{defn}

ఇన్‌పుట్ $w = \sigma_{i_1}\dots\sigma_{i_k}$పై ట్యూరింగ్ యంత్రం~$M$
చర్యను వర్ణించే \tetoken{వాక్యాలు}{!!{sentence}s} ఇవి:
\begin{enumerate}
\item సంఖ్యలనూ~$<$నూ వర్ణించే స్వీకృతాలు:
\begin{enumerate}
\item ప్రతి సంఖ్య దాని ఉత్తరవర్తికంటే తక్కువ అని చెప్పే \tetoken{ఒక వాక్యం}{!!^a{sentence}}:
\[
\lforall[x][x < x']
\]
\item $<$ సంక్రమకమని నిర్ధారించే \tetoken{ఒక వాక్యం}{!!^a{sentence}}:
\[
\lforall[x][\lforall[y][\lforall[z][
      ((x < y \land y < z) \lif x < z)]]]
\]
\end{enumerate}
\item ఇన్‌పుట్ స్థితివిన్యాసాన్ని వర్ణించే స్వీకృతాలు:
\begin{enumerate}
\item $0$~అంచెల తరువాత---యంత్రం మొదలవక ముందు---$M$ ప్రారంభ స్థితి~$q_0$లో
  ఉండి, గడి~$1$ను చదువుతోంది:
\[
\Obj Q_{q_0}(\num{1}, \num{0})
\]
\item మొదటి $k+1$ గడుల్లో $\TMendtape$, $\sigma_{i_1}$, \dots,
  $\sigma_{i_k}$ సంకేతాలు ఉన్నాయి:
\[
\Obj S_\TMendtape(\num{0}, \num{0}) \land
\Obj S_{\sigma_{i_1}}(\num{1}, \num{0}) \land
\dots \land
\Obj S_{\sigma_{i_k}}(\num{k}, \num{0})
\]
\item మిగతా టేపంతా ఖాళీగా ఉంది:
\[
\lforall[x][(\num{k} < x \lif \Obj S_\TMblank(x, \num{0}))]
\]
\end{enumerate}
\item ఒక స్థితివిన్యాసం నుండి తరువాతిదానికి సంక్రమణను వర్ణించే స్వీకృతాలు:

కింది వాటి కోసం, ప్రతి $\sigma \in \Sigma$కు
\[
\lforall[z][
  (((z < x \lor x < z) \land \Obj S_\sigma(z, y))
  \lif \Obj S_\sigma(z, y'))]
\]
రూపంలో ఉన్న అన్ని \tetoken{వాక్యాల}{!!{sentence}s} సంయోగాన్ని $!A(x, y)$గా ఉంచుదాం.
``గడి~$m$ వద్ద తప్ప, $n+1$ అంచెల తరువాత టేపు $n$ అంచెల తరువాత ఉన్నట్లే
ఉంది'' అని వ్యక్తం చేయడానికి $!A(\num{m},\num{n})$ను ఉపయోగిస్తాం.
\begin{enumerate}
\item \ollabel{rep-right} ప్రతి నిర్దేశం $\delta(q_i, \sigma) =
  \tuple{q_j, \sigma', \TMright}$కు, కింది \tetoken{వాక్యం}{!!{sentence}}:
\begin{align*}
& \lforall[x][\lforall[y][(
   (\Obj Q_{q_i}(x, y) \land \Obj S_{\sigma}(x, y)) \lif {}]] \\
&\qquad   (\Obj Q_{q_j}(x', y') \land \Obj S_{\sigma'}(x, y') \land
!A(x, y)))
\end{align*}
ఇది కిందిది చెబుతుంది: $y$ అంచెల తరువాత యంత్రం స్థితి~$q_i$లో ఉండి,
సంకేతం~$\sigma$ గల గడి~$x$ను చదువుతుంటే, $y+1$ అంచెల తరువాత అది
గడి~$x+1$ను చదువుతూ, స్థితి~$q_j$లో ఉంటుంది; గడి~$x$లో ఇప్పుడు~$\sigma'$
ఉంటుంది; $x$ కాని ప్రతి గడిలో $y$ అంచెల తరువాత ఉన్న సంకేతమే ఉంటుంది.

\item \ollabel{rep-left} ప్రతి నిర్దేశం $\delta(q_i, \sigma) =
  \tuple{q_j, \sigma', \TMleft}$కు, కింది \tetoken{వాక్యం}{!!{sentence}}:
\begin{align*}
& \lforall[x][\lforall[y][
    ((\Obj Q_{q_i}(x', y) \land \Obj S_{\sigma}(x', y)) \lif {}]]\\
& \qquad   (\Obj Q_{q_j}(x, y') \land \Obj S_{\sigma'}(x', y') \land
!A(x', y))) \land {}\\
& \lforall[y][((\Obj Q_{q_i}(\num{0}, y) \land \Obj S_{\sigma}(\num{0},
    y)) \lif {}]\\
& \qquad (\Obj Q_{q_j}(\num{0}, y') \land \Obj S_{\sigma'}(\num{0},
  y') \land !A(\num{0}, y)))
\end{align*}
ఇది ఎలా పనిచేస్తుందో కాసేపు ఆలోచించండి: ఇప్పుడు ``గడి~$x$ను చదువుతుంటే
\dots'' అని కాక, ``గడి $x+1$ను చదువుతుంటే \dots'' అని మొదలుపెడతాం. ఎడమకు
కదలడం అంటే తరువాతి అంచెలో యంత్రం గడి~$x$ను చదువుతుంది. కానీ వ్రాయబడే గడి
$x+1$. వ్యవకలనం లేదా పూర్వవర్తి ప్రమేయం మనకు లేవు కనుక ఇలా చేస్తాం.

$x+1$ రూపంలోని సంఖ్యలు $1$, $2$, \dots{} అని గమనించండి. అందువల్ల యంత్రం
గడి~$0$ను చదువుతూ ఎడమకు కదలవలసిన సందర్భం దీనిలో లేదు (అది కదలలేదు;
అక్కడే ఉంటుంది). ఆ ప్రత్యేక సందర్భాన్ని రెండవ సంయోగం చూసుకుంటుంది: $y$
అంచెల తరువాత యంత్రం స్థితి~$q_i$లో ఉండి గడి~$0$ను చదువుతూ, గడి~$0$లో
సంకేతం~$\sigma$ ఉంటే, $y+1$ అంచెల తరువాత అది ఇంకా గడి~$0$నే చదువుతూ,
ఇప్పుడు స్థితి~$q_j$లో ఉంటుందని, గడి~$0$పై సంకేతం~$\sigma'$ ఉంటుందని,
గడి~$0$ కాని గడుల్లో $y$ అంచెల తరువాత ఉన్న సంకేతాలే ఉంటాయని అది చెబుతుంది.\sourcecorrection{OLTETURUND-006}{ఎడమకు కదిలే మొదటి సంక్రమణలో ఆంగ్ల మూలం వ్రాయబడే గడి $x'$ అయినప్పటికీ $!A(x,y)$తో గడి~$x$ను మార్పు-మినహాయింపుగా పెట్టింది. నిర్వచనానికీ వెంటనున్న వివరణకూ అనుగుణంగా, వ్రాయబడే గడి~$x'$ తప్ప మిగతా టేపు మారదని చెప్పే $!A(x',y)$ను ఉపయోగించాం.}
\item \ollabel{rep-stay} ప్రతి నిర్దేశం $\delta(q_i, \sigma) =
  \tuple{q_j, \sigma', \TMstay}$కు, కింది \tetoken{వాక్యం}{!!{sentence}}:
\begin{align*}
& \lforall[x][\lforall[y][(
   (\Obj Q_{q_i}(x, y) \land \Obj S_{\sigma}(x, y)) \lif {}]] \\
&\qquad   (\Obj Q_{q_j}(x, y') \land \Obj S_{\sigma'}(x, y') \land
!A(x, y)))
\end{align*}
\end{enumerate}
\end{enumerate}
పై \tetoken{వాక్యాలన్నిటి}{!!{sentence}s} సంయోగాన్ని, ట్యూరింగ్ యంత్రం~$M$, ఇన్‌పుట్~$w$లకు
$!T(M, w)$గా ఉంచుదాం.

$M$ చివరకు ఆగుతుందని వ్యక్తం చేయడానికి, ``కొన్ని అంచెల తరువాత సంక్రమణ
ప్రమేయం నిర్వచితం కాదు'' అని చెప్పే \tetoken{వాక్యాన్ని}{!!{sentence}} కనుగొనాలి.
$\delta(q, \sigma)$ నిర్వచితం కాని అన్ని యుగ్మాలు $\tuple{q, \sigma}$ల
సమితిని $X$గా ఉంచుదాం. అప్పుడు $!E(M, w)$ను కింది \tetoken{వాక్యంగా}{!!{sentence}} ఉంచుదాం:
\[
\lexists[x][\lexists[y][(\bigvee_{\tuple{q, \sigma} \in
      X}(\Obj Q_q(x, y) \land \Obj S_\sigma(x, y)))]]
\]

ప్రత్యేకంగా నిర్ణయించిన ఆగే స్థితి~$h$ గల ట్యూరింగ్ యంత్రాన్ని ఉపయోగిస్తే
ఇది మరింత సులభం: అప్పుడు కింది \tetoken{వాక్యం}{!!{sentence}}~$!E(M, w)$
\[
\lexists[x][\lexists[y][\Obj Q_h(x, y)]]
\]
యంత్రం చివరకు ఆగుతుందని వ్యక్తం చేస్తుంది.

\begin{prop}
\ollabel{prop:mlessk}
$m < k$ అయితే, $!T(M, w) \Entails \num{m} < \num{k}$.
\end{prop}

\begin{proof}
అభ్యాసం.
\end{proof}

\begin{prob}
\olref[tur][und][rep]{prop:mlessk}ను నిరూపించండి.
(సూచన: $k-m$పై ఆగమనాన్ని ఉపయోగించండి.)
\end{prob}

\end{document}
